\section{Introduction and Dataset}
\label{sec:introduction}

\subsection{Problem Statement and Business Understanding}
\label{sec:problem}

Modern cities depend on a distribution network that stays inside its safe
operating range while demand moves with the weather and with daily activity.
In Tetouan the measured load changes through the day and across the year.
A forecast that arrives after the peak has already formed does not help the
operator. The useful forecast is the one issued before the peak, with enough
lead time to commit reserve capacity.

The target of this project is a short-term forecast of city electricity
demand. The primary horizon is the next hour. A 24-hour horizon is a natural
extension and is left as future work in Section~\ref{sec:conclusion}. The
forecast supports two operating decisions:

\begin{enumerate}[label=(\arabic*), itemsep=2pt]
  \item \textbf{Reserve-margin planning.} Compare the forecast with recent
        history and decide how much spare capacity to hold.
  \item \textbf{Peak-capacity alerting.} Raise an alert when the forecast
        exceeds a capacity threshold, so operators can act one to two hours
        before the highest-risk evening window.
\end{enumerate}

The capacity threshold \(\tau\) used in this study is the 95th percentile of
hourly demand on the training block. It is an empirical stand-in for a
published equipment rating. Replacing \(\tau\) with a utility-defined limit
does not change the rest of the pipeline.

\subsection{Primary Dataset Description}
\label{sec:dataset}

The dataset is the open UCI collection \emph{Power Consumption of Tetouan
City} \cite{uci-tetouan, salam2018}. It records 10-minute observations for
calendar year 2017. Each row contains weather measurements and the power of
three distribution zones. Table~\ref{tab:columns} lists the fields used in
this project.

\begin{table}[h]
  \centering
  \small
  \begin{tabular}{@{}llp{7.2cm}@{}}
    \toprule
    Field & Role & Description \\
    \midrule
    DateTime & Index & Timestamp of the 10-minute sample \\
    Temperature & Weather & Ambient temperature (\(^\circ\mathrm{C}\)) \\
    Humidity & Weather & Relative humidity \\
    Wind Speed & Weather & Wind speed \\
    General diffuse flows & Weather & General diffuse solar radiation \\
    Diffuse flows & Weather & Diffuse solar radiation \\
    Zone 1, Zone 2, Zone 3 & Load & Power of the three distribution zones (kW) \\
    \bottomrule
  \end{tabular}
  \caption{Fields retained from the UCI Tetouan series.}
  \label{tab:columns}
\end{table}

City demand at a 10-minute stamp is the sum of the three zones. Forecasting
is performed on an hourly series. Inside each clock hour the pipeline averages
the weather variables and averages city demand. Averaging keeps the target in
kilowatts of mean power. Summing the six raw power readings inside an hour
would inflate the unit and would no longer match a capacity rating expressed
in kW.

After aggregation, the hourly table is the modelling table. The first 168
hours are consumed by the weekly lag, so the supervised sample starts once
\texttt{lag\_168} is defined. Exact row counts, the calendar span, and the
mean load of each zone are recorded in \texttt{outputs/metrics.json} under
\texttt{raw} and \texttt{profile} when \texttt{run\_pipeline.py} is executed.
Fill Table~\ref{tab:dataset-size} from that file.

\begin{table}[h]
  \centering
  \small
  \begin{tabular}{@{}ll@{}}
    \toprule
    Quantity & Value \\
    \midrule
    Raw rows (10-minute) & [XX] \\
    Hourly rows & [XX] \\
    Supervised rows after lag warm-up & [XX] \\
    First timestamp & [YYYY-MM-DD HH:MM] \\
    Last timestamp & [YYYY-MM-DD HH:MM] \\
    Mean city demand & [XX] kW \\
    \bottomrule
  \end{tabular}
  \caption{Dataset size after download and hourly aggregation.}
  \label{tab:dataset-size}
\end{table}
